% !TeX root = ../../Book3.tex
% 纲要标题：### E.1 微分结构与基本对象

\section{微分结构与基本对象}
\label{b3:sec:E01}

普通多项式关系只约束元素本身，微分方程还把元素与反复求导后的值联系起来。把导子作为环的指定结构后，可以对这些关系作理想运算；常数域也随所选导子而确定。

% 纲要细目（与计划纲要同步）：
% * 指定导子、Leibniz 法则、常数、微分理想、商环及分式域上的延拓；分式域延拓适用于任意整环
% * \(K\{y_1,\ldots,y_n\}=K[y_i^{(j)}:j\ge0]\)、微分阶及微分域扩张；指数函数作为微分代数元的例子
% * 区别指定导子的微分代数、Kähler 微分和 DG 代数，比较值域、Leibniz 法则与平方零条件
% #### E.1.1 导子、常数、微分理想与分式域延拓
% #### E.1.2 微分多项式、微分阶与微分域扩张
% #### E.1.3 微分代数、Kähler 微分与 DG 代数

\subsection{导子、常数、微分理想与分式域延拓}
\label{b3:subsec:E01:01}

除另有说明外，本附录的域均为特征零域，且只指定一个导子 \(\delta\)。例如方程 \(\delta(y)=y\) 同时约束元素与它的导数；取商时，也必须使这个指定导子在商环上有定义。

\begin{definition}\label{b3:def:differential-ring}
设 \(A\) 为交换含幺环，\(\delta:A\to A\) 为加法映射。若对任意 \(a,b\in A\) 都满足
\[
\delta(ab)=a\cdot\delta(b)+\delta(a)\cdot b.
\]
则称 \(\delta\) 为 \(A\) 的\term{导子}[derivation]，称 \((A,\delta)\) 为\term{微分环}[differential ring]。若 \(A\) 是域，则称 \((A,\delta)\) 为微分域。称 \(C_A=\ker\delta\) 为常数环；若 \(A\) 的理想 \(I\) 满足 \(\delta(I)\subseteq I\)，则称 \(I\) 为\term{微分理想}[differential ideal]。
\end{definition}

Leibniz 法则给出 \(\delta(1)=0\)，也给出 \(C_A\) 对加法与乘法封闭。在微分域中，\(\delta(a^{-1})=-a^{-2}\cdot\delta(a)\)，所以非零常数的逆仍为常数，\(C_A\) 是域。微分同态是与导子交换的环同态；其核为微分理想。反之，\(\delta(\bar a)=\overline{\delta(a)}\) 在 \(A/I\) 上良定义，恰好要求 \(I\) 是微分理想。

\begin{proposition}\label{b3:prop:derivation-fraction-extension}
设 \(A\) 为整环，\(\delta:A\to A\) 为导子。它唯一延拓到 \(K=\operatorname{Frac}(A)\)，且对任意 \(a,b\in A\)、\(b\ne0\)，有
\[
\delta(a/b)=\frac{b\cdot\delta(a)-a\cdot\delta(b)}{b^2}.
\]
\end{proposition}
\begin{proof}
在域中对 \(bb^{-1}=1\) 求导强制给出上述公式，故至多有一个延拓。若 \(a/b=c/d\)，则 \(ad=bc\)。对此等式求导，再用 \(ad=bc\) 整理，得到
\[
d^2\bigl(b\cdot\delta(a)-a\cdot\delta(b)\bigr)
=b^2\bigl(d\cdot\delta(c)-c\cdot\delta(d)\bigr).
\]
故公式不依赖代表元。通分直接验证加法与 Leibniz 法则，于是确实得到导子。
\end{proof}

例如在 \(\mathbb Q(t)\) 上指定 \(\delta(t)=1\)，即通常的形式求导。若改为零导子，同一个底层域便成为另一微分域，其常数域也随之改变。因此仅给出一个环，还没有指定微分结构。

\subsection{微分多项式、微分阶与微分域扩张}
\label{b3:subsec:E01:02}

\begin{definition}\label{b3:def:differential-polynomial-ring}
设 \((K,\delta)\) 为特征零微分域，\(n\ge1\)。取彼此代数独立的符号 \(y_i^{(j)}\)（\(1\le i\le n,j\ge0\)），在多项式环
\[
K\{y_1,\ldots,y_n\}=K[y_i^{(j)}:1\le i\le n,\ j\ge0]
\]
上把导子延拓为 \(\delta(y_i^{(j)})=y_i^{(j+1)}\)。所得环称为\term{微分多项式环}[differential polynomial ring]。记 \(y_i^{(0)}=y_i\)。非恒定微分多项式的\term{微分阶}[differential order]是其中实际出现的最高导数阶数。
\end{definition}

每个微分多项式仍只有有限项。例如 \((y')^2-y\) 为一阶、通常多项式次数二，而 \(y^{(3)}-t y\) 为三阶、关于全部导数变量的次数一。次数与微分阶分别记录非线性程度和所用导数深度。

\begin{proposition}\label{b3:prop:differential-polynomial-universal}
设 \((K,\delta)\) 为特征零微分域，\(L/K\) 为保持导子的微分域扩张，\(a_1,\ldots,a_n\in L\)。存在唯一微分 \(K\) 代数同态
\[
\operatorname{ev}_a:K\{y_1,\ldots,y_n\}\longrightarrow L,
\qquad y_i^{(j)}\longmapsto\delta^j(a_i).
\]
\end{proposition}
\begin{proof}
多项式环的泛性质给出唯一普通代数同态。它在系数与所有生成元上和导子交换，再由加法及 Leibniz 法则知在每个多项式上交换。反过来，任一微分同态的所有生成元像都已被 \(y_i\mapsto a_i\) 强制决定。
\end{proof}

设 \((K,\delta)\) 为特征零微分域，\(L/K\) 为保持导子的微分域扩张，\(a\in L\)。若存在非零 \(P\in K\{y\}\) 使 \(P(a)=0\)，则称 \(a\) 为 \(K\) 上的\term{微分代数元}[differentially algebraic element]；否则称为 \(K\) 上的微分超越元。这里 \(P(a)\) 把 \(y^{(j)}\) 替换为 \(\delta^j(a)\)。在带通常求导的亚纯函数域中，\(e^t\) 满足 \(y'-y=0\)，所以是 \(\mathbb C(t)\) 上的微分代数元。普通代数性要求零阶关系，微分代数性则允许导数参与。

\ProseParagraphBreak
记 \(K\langle a\rangle=K(a,\delta a,\delta^2a,\ldots)\)，这是包含 \(K,a\) 的最小微分子域。若 \(a\) 满足关系，并不表示其任意若干初值都能实现为解析解；代数扩张与解析存在性须分别讨论。

\subsection{微分代数、Kähler 微分与 DG 代数}
\label{b3:subsec:E01:03}

设 \(k\) 为域，\(A\) 为交换 \(k\) 代数。本卷第14章的 Kähler 微分 \(\mathrm{d}:A\to\Omega_{A/k}\) 同时表示所有 \(k\) 导子，而这里的 \(\delta:A\to A\) 是一个已经指定的导子。如果 \(\delta|_k=0\)，泛性质给出唯一 \(A\) 线性映射 \(u:\Omega_{A/k}\to A\) 使 \(\delta=u\circ \mathrm{d}\)。因此微分模并不自动选出 \(\delta\)。例如 \(A=k[x]\) 时，\(\Omega_{A/k}=A\,\mathrm{d}x\)，选定 \(u(\mathrm{d}x)=h(x)\) 才得到导子
\[
\delta=h(x)\frac{\mathrm{d}}{\mathrm{d}x}.
\]

另一个不同对象是\term{微分分次代数}[differential graded algebra, DG algebra]。设 \(B=\bigoplus_iB^i\) 为分次 \(k\) 代数，\(\mathrm{d}:B^i\to B^{i+1}\) 为次数一的 \(k\) 线性映射。若 \(\mathrm{d}^2=0\)，且对任意齐次元素 \(a,b\in B\) 满足
\[
\mathrm{d}(ab)=\mathrm{d}(a)b+(-1)^{\deg a}a\,\mathrm{d}(b),
\]
则称 \((B,\mathrm{d})\) 为微分分次代数。本附录的普通导子不要求平方为零。例如在特征零域 \(K\) 的多项式环 \(K[t]\) 中，\(\delta=\mathrm{d}/\mathrm{d}t\) 满足 \(\delta^2(t^2)=2\ne0\)。\cref{b3:tab:E-three-differentials}比较了三种结构：指定导子选择一种求导运算，Kähler 微分表示全部导子，DG 微分则还带有分次和平方零条件。

\begin{table}[htbp]
\centering
\caption{指定导子、Kähler 微分与 DG 微分的比较}
\label{b3:tab:E-three-differentials}
\begin{tabularx}{\linewidth}{@{}>{\raggedright\arraybackslash}p{0.20\linewidth}>{\raggedright\arraybackslash}p{0.28\linewidth}>{\raggedright\arraybackslash}X>{\raggedright\arraybackslash}p{0.17\linewidth}@{}}
\toprule
结构 & 映射与值域 & Leibniz 法则 & 平方零条件 \\
\midrule
指定导子 & \(\delta:A\to A\) & 普通乘积法则 & 不要求 \(\delta^2=0\) \\
Kähler 微分 & \(\mathrm{d}:A\to\Omega_{A/k}\) & 普通乘积法则 & 此映射不能自复合 \\
DG 微分 & \(\mathrm{d}:B^i\to B^{i+1}\) & 含 \((-1)^{\deg a}\) 的分次乘积法则 & 要求 \(\mathrm{d}^2=0\) \\
\bottomrule
\end{tabularx}
\end{table}

多导子版本指定有限个 \(\delta_1,\ldots,\delta_m\)，并要求它们两两交换；生成元变成 \(\delta_1^{\alpha_1}\cdots\delta_m^{\alpha_m}y_i\)。交换性对应混合偏导次序相同。普通微分情形只有一个导子，偏微分情形则要同时处理各个多重导数之间的相容性。基本定义及其计算记号可参阅\cite[Section 3]{BoulierLazardOllivierPetitot2009DifferentialIdeals}。
